\clearpage\section{OBJETIVOS DO CURSO}
\label{OBJETIVOS}

\subsection{Objetivo Geral}

O curso tem como proposta possibilitar a formação de um profissional em engenharia capaz de dominar todas as etapas do desenvolvimento de sistemas de controle e automação de processos e manufaturas, bem como aplicar padrões de engenharia para especificação, dimensionamento e desenho funcional de dispositivos de controle automático de sistemas e unidades de produção. Ao lado da formação técnico-científica, enseja-se a composição de uma perspectiva humanística e empreendedora, criativa e inovadora, crítica e solucionadora de problemas, dando importância ao valor humano, à qualidade de vida e à preservação do meio ambiente.

\subsection{Objetivos Específicos}

O profissional \NOMEEGRESSO possui competências e habilidades para o exercício do cargo conforme as ações previstas na Classificação Brasileira de Ocupações (CBO), com atribuições tais como: planejar serviços, implementar atividades, administrar, gerenciar recursos, promover mudanças tecnológicas e aprimorar condições de segurança, qualidade, saúde e meio ambiente distribuídas nas funções que lhe compete.

Os objetivos específicos do \NOMECURSO são:

\begin{itemize}
	\item Disseminar os conhecimentos sobre aplicações de novas tecnologias com enfoque no controle de processos e na automação industrial;
	\item Viabilizar o trabalho em equipes multidisciplinares, possuindo larga base científica e capacidade de comunicação;
	\item Oportunizar atividades de pesquisa e extensão, estas últimas voltadas preferencialmente às demandas locais e regionais, que favoreçam o desenvolvimento de conhecimento científico e
	tecnológico;
	\item Favorecer a produção de trabalhos científicos, por meio de publicações de alcance local, regional, nacional e
	internacional, com base nos resultados dos trabalhos de conclusão de curso (TCC) e iniciação científica;
	\item Contribuir na inserção dos estudantes no mercado de trabalho de acordo com os arranjos produtivos locais e regionais;
	\item Promover ações para compreensão e aplicação de normas técnicas em saúde, meio ambiente e segurança no trabalho com
	relação às atividades de controle de processos e automação industrial;
	\item Implementar atividades para o desenvolvimento de cultura empreendedora e relações interpessoais;
	\item Avaliar os impactos sociais e ambientais, em especial a nível local e regional, das intervenções inerentes ao cargo e manter o comportamento ético adequado à profissão;
	\item Proporcionar ao graduando uma formação ampla, diversificada, ética e sólida no que se refere aos conhecimentos necessários para a prática profissional;
	\item Promover, por meio das atividades práticas e dos estágios curriculares vivenciados em diversos ambientes de aprendizagem, preferencialmente em empresas localizadas no parque industrial Cearense, a articulação entre teoria e prática;
	\item Contribuir com a inserção dos estudantes em ambientes de produção e divulgação científicas e culturais;
	\item Formar um engenheiro consciente de seu papel no mundo do trabalho nas perspectivas científica, ambiental, ética e social;
	\item Capacitar os futuros engenheiros para assumir a postura de permanente busca de atualização profissional, com vistas a que estes possam compreender, implementar e desenvolver as novas práticas que venham a surgir no campo dos sistemas de controle e automação de processos e manufaturas.
\end{itemize}